% ===========================================================================
% Poste 5 — RAMQ
% ===========================================================================
\poste{5}{Prime au régime public d'assurance médicaments (RAMQ)}{QC\_ramq}

\pdesc Prime annuelle du régime public d'assurance médicaments, payée à la
déclaration de revenus (annexe~K) par les adultes qui ne sont pas couverts
par une assurance collective privée. Le modèle présume la couverture
publique pour tous les adultes.

\pobj Financer l'accès aux médicaments d'ordonnance des personnes sans
régime privé, selon la capacité de payer du ménage.

\pregle La prime d'un adulte dépend du \emph{revenu familial net} du ménage
(équation~\ref{eq:rfn}). On calcule d'abord la base cotisable, c'est-à-dire
l'excédent du $\RFN$ sur le seuil d'exonération $X$, lequel varie selon le
nombre d'adultes et d'enfants :
\[
b = \pos{\RFN - X(\text{adultes},\, \text{enfants})}
\]
puis la prime par adulte, à deux taux (un sur les premiers \D{5000} de base,
un sur l'excédent), plafonnée à la prime maximale $P$ :
\[
\text{prime par adulte} = \min\bigl(\tau_1 \times \min(5\,000,\, b)
 + \tau_2 \times \pos{b - 5\,000},\; P\bigr)
\]
Les couples utilisent un barème à taux réduits de moitié environ, mais
\emph{chacun des deux adultes} paie la prime : le ménage paie donc
$2 \times$ la prime par adulte. Le total est arrondi au cent.

\emph{Exonérations individuelles.} Un adulte ne paie aucune prime s'il
remplit l'une des conditions suivantes, vérifiées par l'orchestrateur
(poste~20) :
\begin{itemize}
  \item il reçoit une aide financière de dernier recours (poste~13) ;
  \item il a 65 ans ou plus et reçoit un Supplément de revenu garanti au
        moins égal à \pc{94} du SRG maximal de sa situation (poste~17).
\end{itemize}

\pparams Le barème (taux et prime maximale) change le 1\ier{} juillet ; les
seuils d'exonération sont indexés au 1\ier{} janvier. Le modèle applique,
pour l'année civile $Y$, le barème de la période tarifaire en cours au
1\ier{}~janvier (débutée en juillet $Y{-}1$) et les seuils indexés de
l'année $Y$ --- c'est la convention du calculateur officiel.

\begin{center}
\small
\begin{tabular}{l r r}
\toprule
Paramètre & 2025 & 2026 \\
\midrule
Prime maximale par adulte ($P$)            & \D{744}   & \D{766} \\
Taux, 1 adulte --- 1\iere{} tranche ($\tau_1$) & \pc{7.65} & \pc{7.84} \\
Taux, 1 adulte --- 2\ieme{} tranche ($\tau_2$) & \pc{11.48} & \pc{11.76} \\
Taux, couple --- 1\iere{} tranche          & \pc{3.84} & \pc{3.93} \\
Taux, couple --- 2\ieme{} tranche          & \pc{5.75} & \pc{5.89} \\
\midrule
Exonération, 1 adulte --- 0 enfant         & \D{19890} & \D{20290} \\
Exonération, 1 adulte --- 1 enfant         & \D{32240} & \D{32890} \\
Exonération, 1 adulte --- 2 enfants et +   & \D{36460} & \D{37195} \\
Exonération, couple --- 0 enfant           & \D{32240} & \D{32890} \\
Exonération, couple --- 1 enfant           & \D{36460} & \D{37195} \\
Exonération, couple --- 2 enfants et +     & \D{40360} & \D{41175} \\
\bottomrule
\end{tabular}
\end{center}

\pnotes Les seuils 2026 sont les seuils indexés utilisés par le calculateur
officiel ; l'annexe~K 2026 n'étant pas encore publiée au moment de la
vérification, ils sont marqués \og à reconfirmer \fg{} dans le code source.

% ============================================================================
% GÉNÉRÉ par scripts/exemples-guide.ts — NE PAS ÉDITER À LA MAIN.
% Régénérer : npm run exemples (chaque chiffre est vérifié contre le moteur).
% ============================================================================
\pexemples La prime se calcule sur le \emph{revenu familial net} ($\RFN$,
tableau~\ref{tab:bases-exemples}), jamais sur le revenu brut.

\medskip\noindent\textbf{M3 --- personne vivant seule, 30~ans, revenu de travail de \D{15000}.}
Son $\RFN$ de \num{13985} est \emph{sous} le seuil d'exonération
(1~adulte, 0~enfant : \D{19890}) : base cotisable nulle, prime nulle.

\medskip\noindent\textbf{M4 --- personne vivant seule, 40~ans, revenu de travail de \D{50000}.}
\begin{align*}
b &= \pos{\num{48115} - \num{19890}} = \num{28225} \\
\text{prime} &= \min\bigl(\pc{7.65} \times \num{5000} + \pc{11.48} \times \num{23225},\ \num{744}\bigr) \\
&= \min(\num{382.50} + \num{2666.23},\ \num{744}) = \num{744}
\end{align*}
La somme des deux tranches dépasse largement la prime maximale : prime
\emph{plafonnée} à \D{744}.

\medskip\noindent\textbf{M6 --- famille monoparentale, 35~ans, revenu de travail de \D{35000}, un enfant : 3~ans (garde subventionnée, \D{2000}).}
Avec un enfant, le seuil d'exonération passe à \D{32240} :
\begin{align*}
b &= \pos{\num{33265} - \num{32240}} = \num{1025} \\
\text{prime} &= \pc{7.65} \times \min(\num{5000},\ \num{1025}) = \num{78.41}
\end{align*}
La base reste dans la première tranche : prime \emph{partielle} de
\D{78.41}.

\medskip\noindent\textbf{M13 --- couple de retraités, 68 et 66~ans, pensions privées de \D{12000} et \D{6000}.}
Couple : barème à demi-taux, seuil de \D{32240}, et la prime
calculée sur le revenu familial commun est payée par \emph{chacun} des deux
adultes.
\begin{align*}
b &= \pos{\num{41237.58} - \num{32240}} = \num{8997.58} \\
\text{prime par adulte} &= \pc{3.84} \times \num{5000} + \pc{5.75} \times \num{3997.58} = \num{421.86} \\
\text{ménage} &= 2 \times \num{421.86} = \num{843.72}
\end{align*}
Prime du ménage : \D{843.72}.

\medskip\noindent\textbf{M1 --- personne vivant seule, 25~ans, aucun revenu.}
Ce ménage reçoit l'aide de dernier recours (poste~13) : chaque adulte
prestataire est \emph{exonéré individuellement} — prime nulle, quel que soit
le calcul ci-dessus.


\prefs Loi sur l'assurance médicaments (RLRQ, c.~A-29.01), art.~10, 23 et
24 ; déclaration TP-1, annexe~K et ligne 447 ; Revenu Québec ; RAMQ (tarifs
et exonérations).
